% !TeX root = ../papers/nc-p3.tex
\section{六行一次表示の核射による階数二曲面の排除}
\label{sec:low-rank-surface-exclusion}
\begin{theorem}\label{thm:low-rank-surface-exclusion}
四辺有限の標準有限Frobenius窓で、入力の二十点は被約、反対平方根 $N$ は全二次式を成分に持ち一般可逆とする。
階数一軌跡が有限なら階数二以下の軌跡の次元は高々1である。
\end{theorem}
\begin{proof}[六行四列の制限が残す次数]
前定理で残った平面 $h=0$ があると仮定する。別の階数二平面があれば、両方の外に少なくとも十六印があり、
両平面を二回ずつ除いた四次式がそれらを三重に通る。これは前定理の四次式の上界に反する。
平面 $h=0$ の外の印を $P$、その数を $n\ge18$ とする。前節末尾の六行線型対称行列を $A$、
その行列式の六次曲面を $G$ とする。$h\ne0$ では $\operatorname{rank} A=2+\operatorname{rank} N$、
従って $P$ の全点で階数3である。$h=0$ で二列 $u,v$ は一般点で独立であり、比例するのは高々二点である。
二列が独立なら $A$ の階数は4以上、また平面上の全点で少なくとも2である。
実際 $u$ は平面の線型埋込みであり、非零の成分とその転置を用いる二行小行列が可逆である。
平面全体で階数4以下なら $h^2\mid\det A$、その残余四次式が $P$ を三重に通ってしまう。
以上より $A$ の階数4以下は高々曲線、3以下は有限で、階数5以下は曲面である。

一般の四次元部分空間 $W\subset k^6$ に制限した $A|_W:4\mathcal O(-1)\to6\mathcal O$ を取る。
その階数3以下の退化は有限である。階数4の点では二次元核に $W$ が非自明に交わる条件が余次元1、
階数5の点では核を $W$ が含む条件が余次元2であり、各々入力層の次元以下に一般繊維の次元を抑える
（入射の次元から一般 $W$ でいずれも零次元以下）。階数6では退化しない。
期待余次元なのでThom--Porteous公式による退化長は
$[H^3](1-H)^{-4}=20$。このうち $P$ が少なくとも $n$ を使う。
特に、$G$ の階数5の成分で核射の像が曲面 $Z$ なら $\deg Z\le20-n\le2$。
実際 $\mathbb P(W)$ は余次元二の一般線型空間で、$Z$ の滑らかな一般点で横断的に切り、
各交点は少なくとも一つの新しい階数5の入力点を与える。
\end{proof}
\begin{proof}[次数五以上の成分]
$G$ が次数 $e\ge5$ の既約成分 $D$ を持てば、次数六なので $G$ は被約である。
$D$ 上の一般核の像の閉包を $Z\subset\mathbb P(U)$、そのベクトル包を $U$、$r=\dim U$ とする。
$A|_U$ の $r$ 行小行列式は全て $D$ 上で消え、一つは一般可逆性から非零である。よって $e\le r$。
特に $Z$ は点でない。像が曲面なら非退化整曲面の次数下界 $\deg Z\ge r-2\ge3$ が直前の上界に反する。

曲線像のとき、一般の $z\in Z$ に対する $A(p)z=0$ は入力の一次方程式であり、入力直線を与える。
実際その線型空間は $G$ に入り、$D$ の一般点を含むため $D$ に含まれる。
一般の繊維次元は1であり、$D$ は平面を含まない。$Z$ の正規化 $B$ から直線のGrassmann多様体への射を
全点に延長し、普遍直線族の引戻し $S$ を取る。$S\to D$ は一般核の一意性から双有理で有限である。
正次元繊維なら全母線が同一の $p$ を通り、$A(p)$ は $r\ge5$ 次元の $U$ を消すが、全点で階数2以上である。
$S$ は滑らかな射影直線束で、$f:S\to G$ をこの有限射、$H$ を入力超平面、$L$ を核超平面の引戻しとする。

$M=\operatorname{coker}(6\mathcal O(-1)\xrightarrow{A}6\mathcal O)$ の表示を $\omega_{\mathbb P^3}$
で双対化すると対称性と超曲面の双対性より
$\mathcal Hom_G(M,\omega_G)=M(-3)$。核座標への評価 $\rho:M\to f_*L$ を双対化し
$\rho(-3)$ と合成すると、有限射の双対性により
\[
         f_*(\omega_S\otimes L^{-1})\longrightarrow M(-3)\longrightarrow f_*L(-3H)
\]
を得る。$D$ の一般点では $G$ は被約、$M$ は階数一で $\rho$ は同型だからこの合成は非零である。
両端は同じ関数体上の階数一捩れなし層の押出しで、合成は一般点での乗法と一致し $f_*\mathcal O_S$ 線型である。
よって $2L-3H-K_S$ に非零切断がある。一方一般母線上でその次数は $0-3+2=-1$ であり矛盾する。
これで全成分の次数が4以下となった。
\end{proof}
\begin{proof}[低次因子における重複度の総和]
印 $P$ での重複度の総和を $m$ と書く。平面は $m\le7$、既約二次曲面は印が高々十四個で頂点の寄与を加えて
$m\le15$。既約三次曲面では $m\le24$ である。三重点がなければ特異集合が有限なら四点以下、
特異直線ならその上の印は三つ以下である。三重点があれば平面三次の錐で、特異母線は高々一本、
頂点でのもう一段を加えても被約な印数に対する増分は4以下である。
従って全因子が三次以下なら重複因子も数えて $m(G)\le 8\deg G=48<3n$。

既約四次因子 $F$ があればその $m(F)$ は $n+19$ 以下である。まず四次錐でなければ四つの偏導三次式は独立であり、
二重点以上の印は高々十六個である。入力三次式が全二十印上の全ての値を与えるため、$s$ 個の印で消える三次式の空間は
次元 $20-s$ である。三重点は前定理の平面三直線の議論で共線となり、その印数は高々三個、四重点は錐以外にない。
ゆえに重複度の増分は $16+3$ 以下である。四次錐なら三つの独立な偏導から二重点以上の印は高々十七個である。
底の整平面四次曲線に三重点がなければ三重以上は頂点のみで、増分は $17+1+1$ 以下である。
底に三重点があれば他の特異点はない（結ぶ直線との交数が4を超える）。特異軌跡は母線一本なので
二重以上・三重以上は各三印以下、四重は頂点のみとなる。以上で同じ上界を得る。
残った次数二の因子の寄与は合わせて15以下であり、$m(G)\le n+34<3n$。
しかし $P$ の全てで $G$ は少なくとも三重である。これで最後の平面も排除された。
\end{proof}
